% Reviewer-motivated rebuttal extension.
% WORKING DRAFT: mathematical protocol only, NO new performance observations.
% Do not claim positive results or submit a results table without H100 receipts.
\section{Token-Granularity Diagnostic for Dirichlet-EESD}
\label{app:token-granularity}

\paragraph{Motivation and scope.}
In the original method, the Dirichlet posterior assigns one supervised
coefficient $w(h)$ to the entire correction $a^1$.
This is a limitation when a corrected program contains many unchanged
tokens as well as a few substantive edits: a useful correction need not
make every copied token equally informative to imitate.
We introduce \emph{edit-aware correction supervision} as an
additional, reviewer-motivated diagnostic rather than a replacement
for the submitted four-transition evidence construction.
It does not alter the recorded experimental results or establish
that token granularity explains the model-dependent outcomes.

\paragraph{Edit-aware weighting.}
For each original/corrected program pair $(a^0,a^1)$, tokenize both
with the same pinned tokenizer used for student training.
Align these token sequences and let $m_{h,t}\in\{0,1\}$ mark
target tokens belonging to inserted or replaced spans;
copied spans are unmarked and the end-of-sequence token is not an edit.
For edit intensity $\eta\ge 0$, define
\begin{equation}
g_{h,t}^{\mathrm{edit}}
=\frac{1+\eta\,m_{h,t}}{1+\eta\,L_h^{-1}\sum_{u=1}^{L_h}m_{h,u}}.
\label{eq:token-aware-gain}
\end{equation}
The denominator guarantees $L_h^{-1}\sum_t g_{h,t}^{\mathrm{edit}}=1$.
This equality is enforced on the \emph{actual response-token prefix}
used in each optimizer micro-step, including when the exact
token budget cuts off the final completion.
When no target token is changed, the weighting reduces to
the original uniform token supervision.
The reference-policy KL remains evaluated on all response positions.

We retain the \emph{original} Dirichlet evidence weight
$w(h)=[\bar u(h)-\lambda\sigma_u(h)]_+$
with $v=(1,-1,0,0)$, $\alpha=0.1$, and $\lambda=0.5$:
\begin{equation}
\mathcal L_{\mathrm{EED+edit}}(h)
=\frac{w(h)}{L_h}\sum_{t=1}^{L_h}
g_{h,t}^{\mathrm{edit}}\bigl[-\log\pi_{t+1}
(a^1_t\mid c_{h,t})\bigr]
+\beta D_{\mathrm{KL}}^{(h)}(\pi_t\Vert\pi_{t+1}),
\label{eq:token-aware-objective}
\end{equation}
where $\beta=0.03$ matches the submitted configuration.
No execution labels are introduced at the token level:
edit membership measures source-code change, not correctness.
Neither the lexical evidence estimator nor the Dirichlet
posterior is replaced by Shapley attribution in this diagnostic.

\paragraph{Matched experimental contrasts.}
Table~\ref{tab:token-aware-protocol} specifies a six-arm design,
chosen to distinguish the evidence-mass construction, token
granularity, and the possibility that concentrating token supervision
is sufficient without knowing the edited positions.
For the random-position control, the number of emphasized target
tokens is matched to the edit mask, and its gain is normalized in
the same way as Equation~\ref{eq:token-aware-gain}.
The edit intensity is fixed to $\eta=2$ before viewing
held-out outcomes. Each arm uses the same original correction
bank, initial LoRA checkpoint, training-source order, learning
rate, response-token budget, KL coefficient, and held-out
evaluation protocol; zero-CE-weight corrections remain in the
schedule for their KL contribution.

\begin{table}[t]
\centering
\small
\caption{\textbf{Proposed reviewer-motivated factorial controls.}
The table defines experimental conditions, not measured scores.
Numerical outcomes must come from new matched GPU runs.}
\label{tab:token-aware-protocol}
\setlength{\tabcolsep}{5pt}
\begin{tabular}{lll}
\toprule
Control & Correction CE weight & Token allocation\\
\midrule
Equal-weight + KL & $1$ & uniform\\
Fixed Dirichlet + KL & $w_n(h)$ & uniform\\
Original EED + KL & $w_{E(p)}(h)$ & uniform\\
Fixed + edit & $w_n(h)$ & $g^{\mathrm{edit}}$\\
EED + edit & $w_{E(p)}(h)$ & $g^{\mathrm{edit}}$\\
EED + random & $w_{E(p)}(h)$ & same-count random $g$\\
\bottomrule
\end{tabular}
\end{table}

The evidence-mass contrast is
$\Delta_{\mathrm{mass}}=P(\mathrm{EED,uniform})
-P(\mathrm{Fixed,uniform})$;
the token contrast is
$\Delta_{\mathrm{edit}}=P(\mathrm{EED,edit})
-P(\mathrm{EED,uniform})$;
and the interaction is
$[P(\mathrm{EED,edit})-P(\mathrm{Fixed,edit})]
-[P(\mathrm{EED,uniform})-P(\mathrm{Fixed,uniform})]$,
where $P$ is independent test-problem all-tests Pass@1.
Reporting a positive token contrast without a positive mass
contrast does not establish that adaptive Dirichlet evidence
improves learning. The random-position control tests whether
specifically emphasizing changed tokens matters beyond
non-uniform supervision.

\paragraph{Data eligibility and limitations.}
This experiment requires a genuine original program $a^0$ and its
paired correction $a^1$ for each training instance.
A pure program-synthesis example with no original program does
not define a source-code edit mask and is excluded from the
\emph{eligible experimental domain}, not selectively removed
after inspecting evaluation performance.
Eligibility, missing source pairs, token-edit fractions,
and source/problem-level train--evaluation separation must be
audited and reported before running or interpreting this study.
In particular, token-diff weighting does \emph{not} itself
correct the submitted utility's equal treatment of
PRESERVE and UNRESOLVED outcomes.
A separate correction-final-quality ablation is needed
to address that concern.
% *** INSERT ONLY OBSERVED, SOURCE-SEALED NUMERICAL RESULTS HERE AFTER
% THE SIX-ARM MATCHED GPU TRAINING AND INDEPENDENT EVALUATION. ***
